\documentclass[12pt]{article}
\usepackage[utf8]{inputenc}
\usepackage[T1]{fontenc}
\usepackage{hyperref}
\title{Tikopia Language Resources for K-12 Curriculum Development}
\author{Compiled by Claude Sonnet 5 (claude-sonnet-5) for Andr\'as Kornai}
\date{}
\begin{document}
\maketitle

\section{Introduction}
Tikopia (ISO 639-3 tkp) is a Polynesian Outlier language of the Austronesian family, spoken on Tikopia in the Solomon Islands. It is closely related to Anuta.

\section{Primary Educational Resources}
No primers, textbooks or school programmes were found in the search results.

\section{Linguistic and Reference Materials}
\begin{itemize}
\item Raymond Firth, \emph{Tikopia-English Dictionary / Taranga fakatikopia ma taranga fakainglisi} (Auckland and Oxford, 1985, with Ishmael Tuki and Pa Rangiaco). It is the main reference, with about 615 pages. Catalogue record: \url{https://ci.nii.ac.jp/ncid/BA04294009?l=en}
\item Firth's draft dictionary manuscript (1979--1992), held at SOAS: \url{https://library.soas.ac.uk/Collection/MS_380598}
\item Entries from Firth 1985 in Pollex-Online (comparative Polynesian lexicon): \url{https://pollex.eva.mpg.de/source/26/?page=22}
\item Bibliographic mentions, with no direct URL found: Dodenhoff, \emph{A Tikopia phonology} (1982); Early, \emph{Minor morphemes of the language of Tikopia} (1981, MA thesis, Auckland).
\item Tikopia language, Wikipedia overview: \url{https://en.wikipedia.org/wiki/Tikopia_language}
\end{itemize}

\section{Religious and Community Resources}
None found.

\section{Digital and Audio Resources}
A language-resource listing was returned (\url{https://cocoon.huma-num.fr/exist/crdo/lexvo.xql?code=tkp}). Its contents were not verified. No audio was confirmed.

\section{Curriculum Development Context}
The language is spoken by all in the ethnic community and is not known to be taught in schools. The only substantial resource found is Firth's dictionary. No full modern grammar was found in these results. A curriculum would have to be developed with the community.
\end{document}